\section{Critical points with degree in the denominator}
\label{cta:critical}

Set
\begin{equation}\label{cta:D}
 b_v=\frac{n_v}{2\pi},\qquad D_v=a_v-\ii b_v,\qquad
 t_0=\frac1{16},\qquad L_v=\lfloor t_0|D_v|\rfloor .
\end{equation}
The summed condition implies
\(\sqrt{d_v^2+n_v^2}\ge2^{25/2}\sqrt{d_v}\), so
\(|D_v|>32\) and \(L_v\ge1\). This denominator stays large
also when \(n_v=0\).

For real continuous charges \(0\le\ell_v\le L_v\), the critical
equation of \(\Phi_{\boldsymbol\ell}\) is
\begin{equation}\label{cta:fixedpoint}
 D_v\tau_v=-\ii\ell_v+\calE_v(\boldsymbol\tau).
\end{equation}
After division by \(D_v\), its right side is a contraction
with constant at most \(\varepsilon\) in the radius-\(1/4\)
polydisk, by \eqref{cta:error} and \(a_v\le|D_v|\).
Its image has radius at most \(t_0+\delta/4<1/4\).
It has a unique analytic critical branch
\(\boldsymbol\tau_{\boldsymbol\ell}\). Define
\(\calC(\boldsymbol\ell)=
\Phi_{\boldsymbol\ell}(\boldsymbol\tau_{\boldsymbol\ell})\).

Put \(Q_v=\calE_v/\tau_v\), with its holomorphic removable
extension. The exact equation is
\begin{equation}\label{cta:tau-formula}
 \tau_v=-\frac{\ii\ell_v}{D_v-Q_v}.
\end{equation}
Since \(|Q_v|\le\delta a_v\), define
\[
 \gamma_+=\frac{1+\delta}{(1-\delta)^2}
          =\frac{2550}{2401},\qquad
 \gamma_-=\frac{1-\delta}{(1+\delta)^2}
          =\frac{2450}{2601}.
\]
Taking the imaginary part of \eqref{cta:tau-formula} gives
\begin{equation}\label{cta:imaginary}
 \gamma_-\frac{\ell_va_v}{|D_v|^2}
 \le-\Ima\tau_v
 \le\gamma_+\frac{\ell_va_v}{|D_v|^2}.
\end{equation}
For example the numerator is
\(\ell_v\Rea(D_v-Q_v)\), between
\((1-\delta)\ell_va_v\) and \((1+\delta)\ell_va_v\);
the denominator is bounded by
\((1\pm\delta)^2|D_v|^2\).
Every positive-charge coordinate therefore has negative
imaginary part and is nonzero.

The envelope derivative is
\(\partial_{\ell_v}\Rea\calC=\Ima\tau_v\).
Integrating it one coordinate at a time gives
\begin{equation}\label{cta:charge-order}
 \Rea\calC(\boldsymbol\ell)\le\Rea\calC(\mathbf1)
   -\frac{\gamma_-}{2}\sum_v
       \frac{a_v}{|D_v|^2}(\ell_v^2-1),
 \qquad 1\le\ell_v\le L_v .
\end{equation}
Integrating from \(\mathbf0\) to \(\mathbf1\) similarly gives
\begin{equation}\label{cta:drop}
 F_0-\beta\le\Rea\calC(\mathbf1)\le F_0,\qquad
 \beta=\frac{\gamma_+}{2}\sum_v\frac{a_v}{|D_v|^2}
                         =\pi\gamma_+\mathfrak S.
\end{equation}
This compares each charge directly with the filled first-charge
action; a comparison only with \(F_0\) would not suffice.

The real part of the Hessian on a real tangent is at most
\(-(1-\varepsilon)\diag(a_v)\), by the symmetric weighted
row bound in \eqref{cta:error}. Its Gaussian determinant
is nonzero. Also
\(|\tau_v|\le t_0/(1-\delta)<\pi/2\), so
\(\prod_v\sin\tau_v\ne0\) at every positive-charge saddle.

\section{The entire Fourier filling calculation}\label{cta:Fourier}

\subsection{A translated finite grid, including its endpoint currents}
For an even trigonometric polynomial
\(f(t)=f^\sharp(-2\cos t)\), define the exact trace
\(\mathcal D_{n,\sigma,N}[f]=\phi_\sigma(R_\sigma^n f^\sharp)\).
There is no restriction on the degree of \(f^\sharp\).
The actual crossing formula of
\cite[Proposition~6.1 and equation~(6.4)]{Meridian} gives
\begin{equation}\label{cta:trace}
 \mathcal D_{n,\sigma,N}[f]
 =\frac{\eta N C_{N,n}(\sigma)}{\sin(2\pi\sigma)}
   \sum_{j=0}^{N-1}
     \sin\tau_j\, f(\tau_j)e^{\ii nN\tau_j^2/(4\pi)},
 \quad
 \begin{cases}
 \tau_j=\dfrac{2\pi}{N}(d/2+\sigma-j),\\
 C_{N,n}(\sigma)=q_N^{n(d^2/4-\sigma^2/2)} .
 \end{cases}
\end{equation}
It is initially written off the colliding grids. Its holomorphic
limit is the full trace, including first normals and the odd
simple-node value; it is not a distinct-node approximation.

Fix \(R<R'<1\), with \(R'\ne1/2\), and first put the formal
parameters \(\sigma_v\) on \(|\sigma_v|=R'\).
Their denominators are then bounded away from zero.
The grid kernel
\[
 \mathcal K_\sigma(t)=
 \frac{\ii N e^{\ii Nt}}{e^{\ii Nt}-\eta e^{2\pi\ii\sigma}}
\]
has residue one at each grid point. Integrate it around the
rectangle with sides \(\pm\pi\pm\ii y+2\pi\sigma/N\).
It encloses exactly the \(N\) points in \eqref{cta:trace}.
On its vertical sides,
\[
 \mathcal K_\sigma(\pm\pi+\ii v+2\pi\sigma/N)
           =\frac{\ii N e^{-Nv}}{1+e^{-Nv}}.
\]
Thus no argument-dependent pole separation is assumed.
The upper and lower horizontal expansions are respectively
\[
 -\ii N\sum_{\ell\ge1}
     \eta^\ell e^{-2\pi\ii\ell\sigma}e^{\ii\ell Nt},
 \qquad
 \ii N\left(1+\sum_{\ell\ge1}
     \eta^\ell e^{2\pi\ii\ell\sigma}e^{-\ii\ell Nt}\right).
\]
Each term has an entire integrand, so its translated path can
be moved to the fixed horizontal segment, retaining the short
endpoint connectors. This is not a deformation of a meromorphic
kernel through its moving poles.

The function \(\sin t\,f(t)e^{\ii nNt^2/(4\pi)}\) is odd.
Reflect the upper side by \(t\mapsto-t\), with its actual
boundary orientation. Combining the two coefficients leaves
\[
 c_{\ell,N}(\sigma)=
   \eta^{\ell-1}U_{\ell-1}(\cos2\pi\sigma),\qquad c_{1,N}=1.
\]
The zero-mode integral on the real segment is zero, but moving
it there leaves its endpoint currents. The resulting
nonzero-charge prefactor is \(-N^2C_{N,n}/(\pi\ii)\).
In particular the reflected charges have already been paired;
there will be no later unexamined sum of equal-action signs.

\subsection{Infinite tails with real spectators}
Let \(A_0=1/400\), and use the heights
\begin{equation}\label{cta:height}
 y_v=A_0t_0\,\frac{a_v}{|D_v|}.
\end{equation}
In coordinate \(v\), with all spectators real or on their
physical translated grids, \eqref{cta:global-Gaussian} bounds
the exponent on the lower horizontal side by
\[
 F_0-A_0a_v t^2+|b_vty_v|-\ell y_v+70a_vy_v^2.
 \]
After maximizing \(t\), the positive quadratic coefficient
of \(y_v^2\) is at most
\[
 70a_v+\frac{100b_v^2}{a_v}
                  \le\frac{100|D_v|^2}{a_v}.
 \]
Since \(L_v+1>t_0|D_v|\), every \(\ell\ge L_v+1\) loses at least
\begin{equation}\label{cta:tail-loss}
 \tfrac12 A_0t_0^2a_v\ge\frac1{409600\pi}.
\end{equation}
Further charges carry \(e^{-N(\ell-L_v-1)y_v}\).
The bound
\(|c_{\ell,N}(\sigma)|\le
\ell e^{2\pi R'(\ell-1)}\) is absorbed for sufficiently large
\(N\), since this tree and \(y_v>0\) are fixed.

The original vertical sides have loss
\(\pi^2a_v/400\).
The chirp, finite-charge connector and imaginary Gaussian
costs there are at most
\(\pi A_0t_0a_v\), \(A_0t_0^2a_v\), and
\(70A_0^2t_0^2a_v\). Their sum is less than half that
endpoint loss. The zero-mode current obeys the same estimate.
The translated-grid error \(O(N^{-1})\) changes only the
prefactor.

For completeness, define the finite real-integral functional
\[
 \mathcal P_v[f]=-\frac{N^2C_{N,n_v}(\sigma_v)}{\pi\ii}
   \sum_{\ell=1}^{L_v}c_{\ell,N}(\sigma_v)
   \int_{-\pi}^{\pi}\sin t\,f(t)
            e^{N(\ii b_vt^2/2-\ii\ell t)}\,\dd t .
\]
Use the exact telescoping identity
\begin{equation}\label{cta:telescoping}
 \prod_v\mathcal D_v-\prod_v\mathcal P_v
 =\sum_j\left(\prod_{v<j}\mathcal P_v\right)
        (\mathcal D_j-\mathcal P_j)
        \left(\prod_{v>j}\mathcal D_v\right).
\end{equation}
The left spectators are finite real integrals and the right
spectators are physical grids. Each costs only a polynomial
in \(N\), with constants depending on fixed powers.
Their grid heights are \(O(N^{-1})\), not another
coordinate's slope-dependent contour height. This is why
the loss in \eqref{cta:tail-loss} survives arbitrary relative
degrees and twists.

Thus the exact retained contraction is, up to an error with
that fixed loss, the finite sum
\begin{equation}\label{cta:finite-charges}
 \begin{split}
 &
 \left(\prod_v-\frac{N^2C_{N,n_v}(\sigma_v)}{\pi\ii}\right)
 \sum_{1\le\ell_v\le L_v}\!
     \left(\prod_vc_{\ell_v,N}(\sigma_v)\right)\\
 &\quad\times
     \int_{[-\pi,\pi]^r}B_N(\boldsymbol\tau,s)
          \prod_v\sin\tau_v\,
          e^{N(\ii\sum_v b_v\tau_v^2/2-\ii\sum_v\ell_v\tau_v)}
                     \,\dd\boldsymbol\tau .
 \end{split}
\end{equation}
The \emph{whole difference} is holomorphic in every formal
\(\sigma_v\) on the unit disk. The original side has the
holomorphic full crossing traces; the finite approximation
has polynomial coefficients \(c_{\ell,N}\).
The maximum principle, one parameter at a time, extends the
estimate from the circles to the disks. We can now set
\(\sigma_v=0,s,-s\) as the actual types require. This step
does not claim that each divided endpoint current separately
has a removable singularity.

\subsection{The connecting contours and the selected charge}
Restrict the real integrals in \eqref{cta:finite-charges} to
\([-\rho,\rho]^r\), with \(\rho=1/8\).
Their complement loses at least \(1/(51200\pi)\) by
\eqref{cta:global-Gaussian}; the real chirps and Fourier
factors have modulus one.
Translate this box to the imaginary part of its own critical
point. On a connector write \(\tau_v=t_v-\ii v_v\), where
\(0\le v_v\le-\Ima\tau_{\boldsymbol\ell,v}\).
All these points lie in the radius-\(1/4\) polydisk.
Equation \eqref{cta:error} and Young's inequality give
\begin{equation}\label{cta:connector}
 \begin{split}
 \Rea\Phi_{\boldsymbol\ell}(\mathbf t-\ii\mathbf v)
 \le F_0-\tfrac18\sum_va_vt_v^2
 +\sum_v\left[
  \tfrac{1+\varepsilon}{2}a_v+
       \frac{2b_v^2}{(3-4\varepsilon)a_v}\right]v_v^2
 -\sum_v\ell_vv_v .
 \end{split}
\end{equation}
Indeed Young's inequality uses
\((3-4\varepsilon)a_vt_v^2/8\) to absorb \(b_vt_vv_v\).
By \eqref{cta:imaginary},
\[
 \gamma_+\max\{(1+\varepsilon)/2,\ 2/(3-4\varepsilon)\}<3/4.
\]
The last two sums in \eqref{cta:connector} are therefore at
most \(-\tfrac14\sum_v\ell_vv_v\). Each connector has some
\(|t_v|=\rho\), and loses at least \(1/(1024\pi)\).
The negative Fourier terms of all spectators are retained
in this calculation; replacing them by independent imaginary
growth estimates would lose the result.

On the translated horizontal box the critical point has
\(|\Rea\tau_v|<t_0/(1-\delta)<\rho\), and the real Hessian
is strictly negative. Apply the holomorphic packet expansion
\eqref{cta:packet-expansion} there. Taylor expansion at that
critical point gives its Gaussian contribution; the rest has
an exponential loss, and the odd Gaussian corrections vanish.
For a nonminimal charge the original integral equals this
saddle contribution plus an additive error at rate
\(F_0-\eta\), where
\begin{equation}\label{cta:eta}
 \eta=\frac1{819200\pi}.
\end{equation}
We do not assert a relative asymptotic for its original
contour if its own action is smaller than that error.
Equation \eqref{cta:charge-order} is what compares its saddle
with the first one.

Every discarded source, edge-contour, real-box, tail and
connector error has at least the loss \eqref{cta:eta}, allowing
polynomial prefactors. Finally
\begin{equation}\label{cta:budget}
 \beta\le\pi\frac{2550}{2401}\,2^{-25}
                      <\frac1{1638400\pi}=\frac\eta2.
\end{equation}
All fixed-loss errors are strictly below the first-charge
action. Each other retained positive charge is also strictly
below it by \eqref{cta:charge-order}. There are finitely many
such charges for the fixed tree. These are two distinct
comparisons, and both are needed.
